% Part: turing-machines
% Chapter: undecidability
% Section: halting-problem

\documentclass[../../../include/open-logic-section]{subfiles}

\begin{document}

\olfileid[fr]{tur}{und}{hal}
\olsection{Le problème de l'arrêt}

\begin{explain}
Supposons fixée une énumération des descriptions de machines de
Turing. Chaque machine reçoit ainsi un \emph{indice}, sa position
dans l'énumération $M_0$, $M_1$, $M_2$, \dots{} des descriptions.\footnote{Précision éditoriale : nous conservons ici l'indexation à partir de $0$ adoptée pour la machine universelle, afin que chaque entier naturel désigne une machine. La dénombrabilité des machines s'entend à renommage des états et des symboles près.}

Nous savons qu'il existe des fonctions non calculables par machine
de Turing : l'ensemble des descriptions de telles machines ---
et donc, à renommage près, celui des machines --- est
!!{enumerable}, tandis que l'ensemble de toutes les fonctions
de $\Nat$ vers~$\Nat$ ne l'est pas. Nous pouvons aussi donner
des exemples particuliers de fonctions non calculables.
La fonction de l'arrêt en est un.
\end{explain}

\begin{defn}[Fonction de l'arrêt]
La \emph{fonction de l'arrêt}~$h$ est définie par
\[
h(e,n) =
\begin{cases}
  \text{0} & \text{si la machine~$M_e$ ne s'arrête pas sur l'entrée~$n$} \\
  \text{1} & \text{si la machine~$M_e$ s'arrête sur l'entrée~$n$}
\end{cases}
\]
\end{defn}

\begin{defn}[Problème de l'arrêt]
Le \emph{problème de l'arrêt} consiste à déterminer, pour tous
$e$ et~$n$, si la machine de Turing~$M_e$ s'arrête sur une entrée
de~$n$ symboles $\TMstroke$.
\end{defn}

\begin{explain}
Nous montrerons que~$h$ n'est pas calculable par machine de Turing
en prouvant d'abord la non-calculabilité d'une fonction apparentée,
notée~$s$. La preuve utilise deux faits déjà établis : tout calcul
effectuable par une machine de Turing peut l'être par une machine
disciplinée (\olref[mac][dis]{sec}), et deux machines peuvent
être raccordées pour n'en former qu'une (\olref[mac][cmb]{sec}).
\end{explain}

\begin{defn}
La fonction~$s$ est définie par
\[
s(e) =
\begin{cases}
  \text{0} & \text{si la machine~$M_e$ ne s'arrête pas sur l'entrée~$e$} \\
  \text{1} & \text{si la machine~$M_e$ s'arrête sur l'entrée~$e$}
\end{cases}
\]
\end{defn}

\begin{lem}
La fonction~$s$ n'est pas calculable par machine de Turing.
\end{lem}

\begin{proof}
Supposons au contraire que~$s$ soit calculable par une machine
de Turing~$S$. Nous pouvons supposer sans perte de généralité
que~$S$ est disciplinée : à l'arrêt, sa tête lit la première
case. Raccordons~$S$ à une machine~$J$, sur un alphabet commun, qui s'arrête si elle
reçoit l'entrée~$0$, c'est-à-dire si, dans son état initial,
elle lit~$\TMblank$ dans la case située immédiatement à droite
du marqueur de début. Sur toute autre entrée, $J$ avance sans
fin vers la droite et ne s'arrête jamais. La machine raccordée
$S \concat J$ est une machine de Turing; elle apparaît donc dans
l'énumération sous un certain indice~$e$ : c'est~$M_e$.
Lançons~$M_e$ sur une entrée de~$e$ symboles $\TMstroke$.
Deux cas sont possibles.
\begin{enumerate}
\item Supposons que~$M_e$ s'arrête sur cette entrée. Alors
  $s(e)=1$. Lancée sur~$e$, la machine~$S$ s'arrête donc avec
  un seul~$\TMstroke$ en sortie. La machine~$J$ démarre alors
  sur un ruban portant ce symbole et ne s'arrête pas. Or~$M_e$
  est précisément $S \concat J$ : elle devrait exécuter~$S$,
  puis~$J$, et donc avancer indéfiniment vers la droite. Elle
  ne peut donc pas s'arrêter sur l'entrée de~$e$ symboles
  $\TMstroke$, contrairement à notre supposition.

\item Supposons maintenant que~$M_e$ ne s'arrête pas sur cette
  entrée. Alors $s(e)=0$ et~$S$, lancée sur~$e$, s'arrête avec
  la sortie vide qui représente~$0$. Lisant la case~$1$ blanche, $J$ s'arrête
  immédiatement. La machine~$M_e=S \concat J$ doit donc elle
  aussi s'arrêter sur l'entrée de~$e$ symboles $\TMstroke$,
  ce qui contredit à nouveau notre supposition.
\end{enumerate}
Les deux cas sont impossibles. Aucune machine~$S$ ne calcule
donc~$s$.
\end{proof}

\begin{thm}[Insolubilité du problème de l'arrêt]
\ollabel{thm:halting-problem}
Le problème de l'arrêt est insoluble : la fonction~$h$ n'est
pas calculable par machine de Turing.
\end{thm}

\begin{proof}
Supposons que~$h$ soit calculable par une machine de Turing~$H$.
Nous pourrions construire une machine calculant~$s$ : elle ferait
d'abord une copie de l'entrée, séparée de l'original par un
symbole~$\TMblank$, reviendrait au début du ruban puis
exécuterait~$H$. Une machine peut effectuer la copie
(voir l'\cref{tur:mac:dis:prob:copier}); si~$H$ existait,
nous pourrions la raccorder à cette machine copieuse pour
décider si~$M_e$ s'arrête sur l'entrée~$e$, autrement dit
pour calculer~$s$. Nous avons démontré que c'est impossible.
La fonction~$h$ n'est donc pas calculable par machine de Turing.
\end{proof}

\begin{prob}
Le problème de l'arrêt sur trois traits, ou 3-Halt,
demande une procédure de décision qui détermine si une machine
de Turing arbitraire s'arrête sur une entrée de trois symboles
$\TMstroke$ inscrits sur un ruban autrement blanc. Démontrez
que ce problème est insoluble.
\end{prob}

\begin{prob}
Montrez que si le problème de l'arrêt était soluble pour les
couples machine--entrée $M_e,n$ avec $e\neq n$, il le serait
aussi pour les cas où $e=n$.
\end{prob}

\begin{prob}
Nous avons montré que le problème de l'arrêt est insoluble
lorsqu'une entrée~$e$ désigne une machine~$M_e$ au moyen d'une
énumération. Qu'en est-il si l'on donne directement en entrée
la description d'une machine de Turing, telle qu'elle a été
présentée dans la \olref[tur][und][enu]{sec} ? Une machine de Turing
pourrait-elle décider le problème de l'arrêt en recevant les
descriptions plutôt que les indices ? Expliquez pourquoi.
\end{prob}

\begin{prob}
Montrez que la fonction \emph{partielle}~$s'$ définie par
\[
s'(e) =
\begin{cases}
  \text{1} & \text{si la machine~$M_e$ s'arrête sur l'entrée~$e$}\\
  \text{non définie} & \text{si la machine~$M_e$ ne s'arrête pas sur l'entrée~$e$}
\end{cases}
\]
\emph{est} calculable par machine de Turing.
\end{prob}

\end{document}
